\shortname{Sorting Algorithms}
\style{sorting.css}
\script{../core.js}
\script{sound.js}
\script{generate_array.js}
\script{sorting.js}
\script{simple_sorts.js}
\script{clever_sorts.js}
\script{silly_sorts.js}
\script{compare_sorts.js}

\section{A Needlessly Thorough Analysis of Known Sorting Algorithms}
Introduction

\subsection{Compare sorting algorithms}
\div{compare_sort}{compare_algo}

\subsection{Simple sorting algorithms}
\subsubsection{Selection sort}
\div{selection_sort}{sorting_algo}
The simplest sorting algorithm is selection sort, this algorithm starts by finding
the smallest value in the array and placing it in the front. Afterwards it will find the
smallest value in the array, but this time it will exclude the first value, subsequently 
finding the second smallest value and placing it in the second slot in the array. 
Doing this continously will eventually lead to a sorted array. 
<h4>Pseudocode</h4>
\code{##SelectionSort##(~A~)
    for i = 0 to |~A~| - 2
        //This number will denote the index of the smallest element in the array
        minIndex = i 
        for j = i + 1 to |~A~| - 1
            if ~A~[j] < ~A~[minIndex]
                minIndex = j
            endif
        endfor
        
        //When we've found the smallest element, we place it on the i'th index
        ##Swap##(~A~, i, minIndex) 
    endfor

    return ~A~}

<h4>Time- and spatial complexity</h4>
No matter how scrambled 
the original array was, this algorithm will sort the array in $\frac{n(n-1)}{2}-1\in
\Theta(n^2)$ steps 
(where $n$ is the size of the original array). To see why, note that to find the smallest
value we first have to check $n$ elements, then to find the next smallest value we check
the remaining $n-1$ elements and so on
\[
    (n-1)+(n-2)+\cdots+3+2=\sum_{i=2}^{n-1} i=\sum_{i=1}^{n-1} i-1=\frac{n(n-1)}{2}-1\in\Theta (n^2)
\]
We could also look at it a different way, using upper and lower bounds, let $f(n)$ be the
number of operations performed by SelectionSort for an array of size $n$. The following
diagram indicates the amount of work performed by our algorithm according to the value
of $i$ and $j$. Wether or not the if statement meets is condition is the difference 
between 1 operation and 2, so its up to a constant and can be dismissed.
\\\\
\fig{complexity.png}{70}
\figtext{On the left: number of operations performed by insertion sort per index $i$ and $j$. <br>
On the right: operations performed by insertion sort for an array of size 5.}

Notice that we can fill a red square of size $n^2/4$ where all work inside the 
square is 1. Thus the total work performed in this square is $n^2/4$, this is a
lower bound so $f(n)\in\Omega(n^2)$. Now the entire square of size $(n-1)^2$ is an 
overestimate (combining the green and red squares) since not all the work in this 
large square is 1, thus we also have an upper bound $f(n)\in\mathcal O (n^2)$. Since the 
upper and lower bounds are the same we have our precise estimate $f(n)\in \Theta(n^2)$.
\\\\
Now in terms of spatial complexity we create only a variable <code>minIndex</code>
and a temporary variable for swapping elements, so in total two variables, which
makes the spatial complexity constant $\Theta(1)$.
<h4>Proof of correctness</h4>
Let $\mathbf A =\left< a_0,a_1,\cdots,a_{n-1}\right>$ be a list of $n$ elements.
This method of proof uses an $invariant$. This invariant is the following:
$\forall k< i,\forall l\geq i: a_k\leq a_l$ - in words: all elements to the right of 
$a_{i-1}$ are larger than or equal to all elements to the left of $a_i$.
\\\\
Base case:\\
This holds trivially for $i=0$, since there is no $j< i$.
\\\\
Maintenance:\\
For every loop with respect to $i$, $j$ will go through all elements of $\mathbf A$
to the right of $i-1$ and keep track of the smallest one. If the smallest element
is found to be $a_m$ then $a_i$ and $a_m$ swap places. Recall that $a_m$ was smaller
than or equal to all elements to the right of $i-1$, thus since our invariant holds
for $i-1$ it holds also for $i$.
\\\\
Termination:\\
Once $i=n-1$ we get that the element at $n-1$ is the largest and that all elements to the
left of $n-1$ are in sorted order (since $a_j$ $\forall j< i$ is smaller than any element 
to the right of it - by our invariant)

\subsubsection{Bubble sort}
\div{bubble_sort}{sorting_algo}
<h4>Pseudocode</h4>
\code{##BubbleSort##(~A~)
    for i = 0 to |~A~| - 2
        swapped = false
        for j = 0 to |~A~| - i - 1
            //Swap all elements that are larger than the element after it
            if ~A~[j] > ~A~[j + 1]
                ##Swap##(~A~, j, j + 1) 
                swapped = true
            endif
        endfor

        //If no elements are swapped we can stop early
        if not swapped
            break
        endif
    endfor
        
    return ~A~}

<h4>Time- and spatial complexity</h4>
Bubble sort only uses extra storage when swapping elements, thus the spatial complexity is $\O(1)$.\\
The time complexity for bubble sort is in the worst case and average case $\O(n^2)$, but $\O(n)$ in the 
best case.
\\\\
Because our implementation of \textsc{BubbleSort} allows for a quick exit then the array is sorted the time
complexity gets a bit more complicated to calculate, but not impossible. \\\\
We start our analysis by letting $h(i)$ be a function that for each index in $\A$ returns the index where 
$\A[i]$ belongs in a sorted array.\\
Now if $h(i)< i$ it would take $a=i-h(i)$ iterations of the outer loop before $\A[i]$ is where it belongs,
since an element only moves back once. If on the other hand $h(i)>i$, then it can be carried to where it belongs
after \textsc{BubbleSort} has placed all elements that will go to the right of $h(i)$, i.e after $b=|\A|-1-h(i)$ 
iterations of the outer loop (except if $h(i)=i$, then it is already in its place and $b=0$). 
Therefore the number of iterations of the 
outer loop is exactly $\mathcal S=1+\displaystyle\max_{0\leq i<|\A|} \big\{\max(a,b)\big\}$ 
(plus one, since we always run the loop before we realize that we didn't have to).\\
From this we find that the runtime of \textsc{BubbleSort} is 
\begin{equation*}
    \begin{split}
        \sum_{i=1}^{\mathcal S} (n - i)=\sum_{i=1}^{\mathcal S} n-\sum_{i=1}^{\mathcal S} i\quad(n=|\A|)\\
        =n\mathcal S-\frac{\mathcal S(\mathcal S + 1)}{2}=
        \Theta\left(\mathcal S\left[n-\frac{\mathcal S + 1}{2}\right]\right)
    \end{split}
\end{equation*}
but since $S\leq n$, $\Theta\left(\mathcal S(n-\frac{\mathcal S + 1}{2})\right)=\Theta(\mathcal S n)$.
Thus the problem of determing time complexity is narrowed down to finding an expression of $\mathcal S$ 
in terms of $n$. Now if the array is sorted $h(i)=i$ for all $i$, therefore $a=i-h(i)=0$ and $b=0$ making 
$\mathcal S=1$. This is the best case scenario and the time complexity is $\Theta(n)$. 
\\
\textsc{BubbleSort} enters a worst case scenario if for example the smallest element was placed at the last index,
if this happens $\mathcal S=1+i-h(i)$ with $i=n-1$ and $h(i)=0$ making $\mathcal S=n$ and the time complexity is 
in turn $\Theta(n^2)$.
\\
In the average case $\mathcal S$ will be proportional to $n$, i.e $\mathcal S=\Theta(n)$, since the average
element would be just as likely to be larger than, as it is likely to be smaller than any other element. 
Meaning the average element is the median of the array. Thus the first element would have $n/2$ elements
larger than it, and $b=|\A|-1-h(i)=n-1-n/2=n/2-1=\Theta(n)$. Making the average case time complexity at least
$\Theta(n^2)$.


\subsubsection{Insertion sort}
\div{insertion_sort}{sorting_algo}
<h4>Pseudocode</h4>
\code{##InsertionSort##(~A~)
    for i = 1 to |~A~| - 1
        key = ~A~[i]
        j = i - 1

        //Shift earlier elements forward until all are larger than the key
        while j >= 0 and ~A~[j] > key
            ~A~[j + 1] = ~A~[j]
            j--
        endwhile

        //Place key into the gap we created by shifting all other elements
        ~A~[j + 1] = key
    endfor
        
    return ~A~}
Although for visualization purposes it made more sense to use this alternative formulation
\code{##AlternativeInsertionSort##(~A~)
    for i = 1 to |~A~| - 1
        j = i - 1

        //Shift A[i] back as long as it's smaller than the element before it
        while j >= 0 and ~A~[j] > ~A~[j + 1]
            ##Swap##(~A~, j, j + 1)
            j--
        endwhile
    endfor
        
    return ~A~}

Let $\mathcal S$ be the number of iterations of the inner while loop, then the time complexity of \textsc{InsertionSort}
is $\Theta(\mathcal S n)$ and $\mathcal S$ is exactly the number of times where $\A[j - 1]>\A[j]$ for any $1\leq j<|\A|$
(note that $\mathcal S_i$, i.e the value of $\mathcal S$ at iteration $i$ is not affected by rearranging the earlier 
elements). 
\\\\
If $\mathcal S_i$ is the number of times the while loop runs in the $i$'th iterations of the for loop, then the
run time for \textsc{InsertionSort} is $=\displaystyle\sum_{i=1}^{|\A|-1}(1+\mathcal S_i)$. The average value for 
$1+\mathcal S$ is equal to
\begin{equation*}
    \begin{split}
        \mathcal S=\frac{1}{n}\sum_{i=1}^{|\A|-1} (1+\mathcal S_i)\\
        \iff\\
        n\mathcal S=\sum_{i=1}^{|\A|-1} (1+\mathcal S_i)=\text{run time}
    \end{split}
\end{equation*}
and thus the run time is just $n$ times the average value of $1+S_i$.\\
In the worst case all the elements are in reversed sorted order, and $S_i$ is equal to $i$ (since the $i$'th element
has to move all the way to) and the average value of $S_i=\frac{n}{2}$ making the worst case run time
$\Theta(n^2)$.\\
In the average case the $i$'th element has $i/2$ elements before it, that is larger than itself. Thus 
$S_i=\frac{n}{4}$ and the average time complexity is $\Theta(n^2)$ aswell.\\
And in best case the array is already sorted, thus $S_i=0$ and the run time is $\Theta(n)$.


\subsection{Clever sorting algorithms}
\subsubsection{Merge sort}
\div{merge_sort}{sorting_algo}

<h4>Pseudocode</h4>
\code{##MergeSort##(~A~)
    //Start the recursion:
    ##MergeSort##(~A~, 0, |~A~|)
    return ~A~
    
##MergeSort##(~A~, ~l~, ~r~)
    //If this partition of the array has only 1 element, it's already trivially sorted
    if ~r~ - ~l~ <= 1
        return
    endif

    m = ##Floor##((~l~ + ~r~) / 2)

    //Recursively sort both halves of the array
    ##MergeSort##(~A~, ~l~, m)
    ##MergeSort##(~A~, m, ~r~)

    //And now merge the two sorted halves
    ##Merge##(~A~, ~l~, m, ~r~)
    
##Merge##(~A~, ~l~, ~m~, ~r~)
    //Copy the left and right halves of the array
    leftHalf = ~A~[~l~..~m~]
    rightHalf = ~A~[~m~..~r~]
    
    //These pointers will point to indices in the left and right half respectively
    L = 0
    R = 0
    //While both halves have elements left pick the smallest element in both arrays
    while L < len(leftHalf) and R < len(rightHalf)
        i = l + L + R

        //Since the halves are sorted, the smallest element in each array is the first element
        //This is what makes merge O(n)
        if leftHalf[L] < rightHalf[R]
            ~A~[i] = leftHalf[L]
            L++
        else
            ~A~[i] = rightHalf[R]
            R++
        endif
    endwhile
    
    //If there are no elements left in the right half, copy the remaining elements of the left half
    while L < len(leftHalf)
        i = l + L + R
        ~A~[i] = leftHalf[L]
        L++
    endwhile
    
    //If there are no elements left in the left half, copy the remaining elements of the right half
    while R < len(rightHalf)
        i = l + L + R
        ~A~[i] = rightHalf[R]
        R++
    endwhile}
\subsubsection{Quick sort}
\div{quick_sort}{sorting_algo}
\code{##QuickSort##(~A~)
    ##QuickSort##(~A~, 0, |~A~|)

##QuickSort##(~A~, ~l~, ~r~)
    if(~r~ <= ~l~)
        return

    mid = ##Partition##(~A~, ~l~, ~r~)

    ##QuickSort##(~A~, ~l~, mid - 1)
    ##QuickSort##(~A~, mid + 1, ~r~)

##Partition##(~A~, ~l~, ~r~)
    x = ~A~[r]
    i = ~l~ - 1
    for j = ~l~ to ~r~ - 1
        if ~A~[j] <= x
            i++
            ##Swap##(~A~, i, j)

    ~A~[~r~] = ~A~[i + 1]
    ~A~[i + 1] = x
    
    return i + 1}

\subsubsection{Heap sort}
\div{heap_sort}{sorting_algo}

\subsubsection{Introspective sort (Intro sort)}
\div{intro_sort}{sorting_algo}

\subsubsection{Shell sort}

\subsection{Theoretical limits for comparison-based sorts}
Assume we have an array $\mathbf A=\left< a_0,a_1,\cdots a_{n-1}\right>$ of $n$ distinct elements, there are exactly $n!$ 
different permutations of this array
and only one of these permutations is the sorted array, so our algorithm will have to distinguish between all
$n!$ permutations to find the right one. Comparison based algorithms do this by comparing elements, to ensure
that smaller elements are to the left of larger elements.
\\\\
A sorting algorithm might start by comparing the first two elements: $a_0< a_1$, 
if this inequality holds then the first two elements are in sorted order, if not they are in reverse sorted
order, no matter what happens we can eliminate half of all possibilites. Assume that $a_0 < a_1$
then all permutations containing with $\mathbf A=\left< \cdots, a_1, \cdots, a_0,\cdots\right>$ (that is,
having $a_1$ before $a_0$) can instantly be discarded. This accounts for half of all permutations, and thus 
after just one step we go from $n!$ permutations to check down to $n!/2$ (note that $a_0$ and $a_1$ are
not unique, you could compare any two indices $0\leq i, j < n,i\neq j$, $a_i< a_j$ in the same way we can show
that the argument of halving the permutations also holds whenever the inequality $a_i< a_j$ does not. Since
this is just the case of $j< i$).
\\
Now we compare the next elements $a_1< a_2$ and again let's assume that this inequality holds, then we can
cut away all permutations containing $\mathbf A=\left< \cdots, a_2, \cdots, a_1,\cdots\right>$ which again
halves the number of permutations we have to check: $n!/4$. We can keep going until we know the exact 
permutation that sorts the array (i.e, until $n!/2^k\leq 1$).
\begin{equation}\label{f n geq n fac}
    \begin{split}
        1\geq \frac{n!}{2^{f(n)}}\\
        2^{f(n)}\geq n!\\
        f(n)\geq \log_2 (n!)
    \end{split}
\end{equation}
where $f(n)$ is the number of steps our theoretical sorting algorithm has to perform. If we expand out
$\log_2(n!)$ we have for even values of $n$:
\begin{equation}\label{log2 n factorial}
    \begin{split}
        \log_2(n!)=\log_2(1)+\log_2(2)+\cdots+\underbrace{\log_2([n+1]/2)+\cdots+\log_2(n)}_{n/2\text{ terms}}\\
        \underbrace{\log_2([n+1]/2)+\cdots+\log_2(n)}_{n/2\text{ terms}}\geq 
        \underbrace{\log_2(n/2)+\cdots+\log_2(n/2)}_{n/2\text{ terms}}\\
        =\frac{n}{2}\log_2{(n/2)}=\Theta(n\log n)
    \end{split}
\end{equation}
for odd values of $n$ we recall that this is a lower bound, so use $\log_2(n!)\geq \log_2((n-1)!)=\Theta(n\log n)$.
\\
Now combine (\ref{f n geq n fac}) with (\ref{log2 n factorial}) and 
we've succesfully arrived at the theoretical limit of the time complexity 
of sorting algorithms is $f(n)=\Omega(n\log n)$. By the existance of algorithms like MergeSort
we know that this theoretical limit is also achievable.

\subsection{Silly sorting algorithms}
\subsubsection{Gnome sort}
\div{gnome_sort}{sorting_algo}
